注意, 由于存在从 $H^k(\mathbb{R}_+^m)$ 到 $H^k(\mathbb{R}^m)$ 的延拓算子, 也有一个迹算子, 将 $H^k(\mathbb{R}_+^m)$ 中的函数映到它在 $H^{k-1/2}(\mathbb{R}^{m-1})$ 中的边界值. 利用单位分解论证, 可以将这一结果推广到具有有界边界的区域.

\begin{Theorem}
\label{thm:smooth-domain-trace-and-right-inverse}
设 $k$ 为正整数. 假设 $\Omega$ 为 $C^k$ 类, $\partial\Omega$ 有界. 那么存在有界迹算子 $\Gamma:H^k(\Omega)\to H^{k-1/2}(\partial\Omega)$. 此外, $\Gamma$ 具有有界右逆.
\end{Theorem}

若 $l<k$, 则第 $l$ 阶导数在 $H^{k-l-1/2}(\partial\Omega)$ 中具有迹. 通常用沿 $\partial\Omega$ 法向方向的导数, 来表述涉及高阶导数的迹定理.

\begin{Theorem}
\label{thm:domain-higher-normal-trace}
设 $k,l$ 为满足 $k>l$ 的正整数. 设 $\Omega$ 为 $C^k$ 类, $\partial\Omega$ 有界. 那么存在连续迹算子
\begin{equation}\label{eq:domain-higher-normal-trace-product-map}
T_l:H^k(\Omega)\to\prod_{j=0}^lH^{k-j-1/2}(\partial\Omega)
\end{equation}
使对每个光滑的 $\phi$ 有
\begin{equation}\label{eq:domain-higher-normal-trace-smooth-values}
T\phi=\left(\phi,\dfrac{\partial\phi}{\partial n},\ldots,\dfrac{\partial^l\phi}{\partial n^l}\right)
\end{equation}
算子 $T_l$ 具有有界右逆.
\end{Theorem}

现在可以用边界条件刻画 $H_0^k(\Omega)$.

\begin{Theorem}
\label{thm:zero-boundary-sobolev-characterization}
设 $\Omega$ 为 $C^k$ 类, $\partial\Omega$ 有界. 那么 $H_0^k(\Omega)$ 就是所有在 $\partial\Omega$ 上以迹的意义满足
\begin{equation}\label{eq:zero-boundary-sobolev-normal-trace-characterization}
u=\dfrac{\partial u}{\partial n}=\cdots=\dfrac{\partial^{k-1}u}{\partial n^{k-1}}=0
\end{equation}
的函数 $u\in H^k(\Omega)$ 所组成的集合.
\end{Theorem}

\begingroup
\postdisplaypenalty=10000
\renewcommand{\endProof}{\par\nobreak\hfill$\blacksquare$\par\vspace{0.1cm}}
\begin{Proof}
若 $u\in\mathcal{D}(\Omega)$, 则显然 \eqref{eq:zero-boundary-sobolev-normal-trace-characterization} 成立. 由连续性, \eqref{eq:zero-boundary-sobolev-normal-trace-characterization} 对 $u\in H_0^k(\Omega)$ 也成立. 需要证明逆命题. 利用单位分解和局部坐标变换, 可化为 $\Omega=\mathbb{R}_+^m$ 的情形. 现在设 $k=1$, $u\in H^1(\mathbb{R}_+^m)$ 在迹的意义下满足 $u(\bm{x}^{\prime},0)=0$. 用 $Eu$ 表示 $u$ 的零延拓. 为证明 $Eu\in H^1(\mathbb{R}^m)$, 只需证明 $\partial(Eu)/\partial x_i=E(\partial u/\partial x_i)$. 对 $i<m$ 这显然成立. 对 $i=m$, 对任意 $\phi\in\mathcal{D}(\mathbb{R}^m)$ 有:
\begin{equation}\label{eq:zero-extension-normal-weak-derivative-calculation}
\begin{aligned}
\int_{\mathbb{R}^m}\dfrac{\partial\phi}{\partial x_m}(Eu)\,\mathrm{d}\bm{x}
&=\int_{\mathbb{R}_+^m}\dfrac{\partial\phi}{\partial x_m}u\,\mathrm{d}\bm{x}\\
&=\int_{\mathbb{R}^{m-1}}\phi(\bm{x}^{\prime},0)u(\bm{x}^{\prime},0)\,\mathrm{d}\bm{x}^{\prime}-\int_{\mathbb{R}_+^m}\phi\dfrac{\partial u}{\partial x_m}\,\mathrm{d}\bm{x}\\
&=-\int_{\mathbb{R}^m}\phi E\dfrac{\partial u}{\partial x_m}\,\mathrm{d}\bm{x}.
\end{aligned}
\end{equation}
类似论证适用于更高阶导数. 一旦知道 $Eu\in H^k(\mathbb{R}^m)$, 其余结论便可通过考察序列 $u_n=Eu(\bm{x}-\dfrac{1}{n}\bm{e}_m)$ 得到. 由于 $u_n$ 的支集与 $\partial\mathbb{R}_+^m$ 保持正距离, 容易用试验函数逼近 $u_n$.
\end{Proof}
\endgroup
